\documentclass{report}
\usepackage{amsmath,amssymb}
\setlength{\parindent}{0mm}
\setlength{\parskip}{1em}
\begin{document}
\begin{center}
\rule{6in}{1pt} \
{\large Frank E. Curtis \\
{\bf Infeasibility Detection in Nonlinear Optimization}}

Lehigh University \\ Department of ISE \\ 200 W Packer Ave \\ Bethlehem \\ PA 18015
\\
{\tt frank.e.curtis@gmail.com}\\
Jim Burke\\
Hao Wang\end{center}

Contemporary numerical methods for nonlinear optimization possess strong
global and fast local convergence guarantees for feasible problems under
common assumptions. They also often provide guarantees for (eventually)
detecting if a problem is infeasible, though in such cases there are
typically no guarantees of \emph{fast} local convergence. This is a
critical deficiency as in the optimization of complex systems, one often
finds that nonlinear optimization methods can fail or stall due to minor
constraint incompatibilities. This may \emph{suggest} that the problem is
infeasible, but without an infeasibility certificate, no useful result is
provided to the user. We present a sequential quadratic optimization
(SQO) method that possesses strong global and fast local convergence
guarantees for both feasible and infeasible problem instances.
Theoretical results are presented along with numerical results indicating
the practical advantages of our approach.


\end{document}
